\chapter*{KATA PENGANTAR}
\addcontentsline{toc}{chapter}{KATA PENGANTAR}

Puji syukur penulis panjatkan ke hadirat Tuhan Yang Maha Esa karena atas rahmat dan karunia-Nya, penulis dapat menyelesaikan penulisan Laporan Tugas Akhir yang berjudul ``RAG Terkondisi-Prediksi: Integrasi Prediksi Kadar Glukosa dan Rekomendasi Klinis Berbasis Dokumen untuk Pengelolaan Diabetes''. Laporan ini disusun sebagai salah satu syarat untuk menyelesaikan program Sarjana di Program Studi Sistem dan Teknologi Informasi, Sekolah Teknik Elektro dan Informatika, Institut Teknologi Bandung.

Penyelesaian tugas akhir ini tidak terlepas dari bimbingan, dukungan, dan doa dari berbagai pihak. Oleh karena itu, penulis menyampaikan terima kasih yang sebesar-besarnya kepada:
\begin{enumerate}
    \item Bapak Prof. Dr. Ir. Suhono Harso Supangkat, M.Eng. selaku Dosen Pembimbing I, yang telah memberikan arahan, kritik, dan saran yang berharga sejak perumusan masalah hingga penyusunan laporan ini.
    \item Bapak Ir. Devi Willieam Anggara, S.T., M.Phil., Ph.D., IPP selaku Dosen Pembimbing II, yang telah meluangkan waktu, tenaga, dan pikiran untuk membimbing penulis selama proses penelitian dan penulisan laporan.
    \item Bapak Prof. Ir. Armein Z. R. Langi, M.Sc., Ph.D. dan Bapak Ishak Hilton Pujantoro Tnunay, S.T., Ph.D. selaku dosen penguji, yang telah memberikan masukan untuk menyempurnakan laporan ini.
    \item Bapak Ir. I Gusti Bagus Baskara Nugraha, S.T., M.T., Ph.D. selaku Ketua Program Studi Sistem dan Teknologi Informasi, Bapak Dr. Ir. Albarda, M.T. selaku dosen wali, serta segenap dosen dan tenaga kependidikan Sekolah Teknik Elektro dan Informatika ITB yang telah memberikan bekal ilmu dan dukungan selama masa perkuliahan.
    \item Para dokter dan tenaga kesehatan yang telah bersedia menjadi responden evaluasi ahli, yang masukannya menjadi bagian penting dari evaluasi sistem pada tugas akhir ini.
    \item Kedua orang tua dan keluarga penulis atas doa, dukungan, dan kasih sayang yang tiada henti selama masa studi.
    \item Vanesa Tri Ananda, S.Ked., yang telah memberikan dukungan moril tanpa henti, menjadi teman diskusi dalam perumusan ide penelitian, serta membantu menghubungkan penulis dengan para responden evaluasi ahli.
    \item Sahabat-sahabat SSS (Jo, Billy, Harman, Tepon, Alpin, Jul, dan Miki), Dama, serta rekan-rekan mahasiswa Sistem dan Teknologi Informasi ITB angkatan 2022, yang telah menjadi teman diskusi dan berbagi semangat selama masa perkuliahan maupun pengerjaan tugas akhir.
    \item Semua pihak yang tidak dapat penulis sebutkan satu per satu, yang telah membantu secara langsung maupun tidak langsung dalam penyelesaian tugas akhir ini.
\end{enumerate}

Penulis menyadari bahwa laporan ini masih memiliki kekurangan. Oleh karena itu, kritik dan saran yang membangun sangat penulis harapkan. Semoga laporan ini dapat memberikan manfaat bagi pengembangan ilmu pengetahuan, khususnya pada bidang sistem pendukung keputusan klinis dan pengelolaan diabetes.
\begin{flushright}
    Bandung, \today\\

    Penulis \\
    \vspace{1cm}
    Daffari Adiyatma
\end{flushright}
